%======================================================================
\section{Track PDM --- Part and drawing management}\label{sec:pdm}
%======================================================================

FreeCAD has no PDM: no vault, check-out, revision, lifecycle, where-used or release. Fusion has
versions in its cloud plus the Manage extension (part numbering, lifecycle, change orders, BOM release).
SolidWorks has PDM Standard and Professional (vault, check-in/out, workflows, revisions, BOMs). Searches
for an open-source PDM with a working FreeCAD check-in workflow found only early or narrow projects
(FreePDM, a concept; a small Taack-based PLM; LibrePLM, with a Windows-only connector). This track
therefore \textbf{starts with an evaluation}, as decided, and specifies the data model and lifecycle
independently of where they end up stored.

\begin{epic}{PDM-01}{PDM evaluation \& decision}{M2}{\stDone}{(new)}
Compare four architectures against weighted criteria and record the decision. \textbf{Decided 19 Oct 2026
(\S\ref{sec:pdm-decision}):} a vault in any folder, iCloud Drive first, opened through a \file{.bwvault} ref file;
one writer at a time for now. The folder layout is designed so a sync service has nothing to conflict over.
\end{epic}

%======================================================================
\subsection{BW-PDM-01 --- Decision: a vault in a folder (19 Oct 2026)}\label{sec:pdm-decision}
%======================================================================

\textbf{Decision (your answer, 19 Oct).} The vault lives in any folder you choose and builds itself there; a ref
file lets a reader view and drive it; the backend is that folder, most likely in iCloud, so it moves easily to
another provider. Scope, also 19 Oct: \textbf{one FreeCAD at a time}; simultaneous edits from several instances are
not part of this step.

\paragraph{Layout.} A sync service makes conflict copies only when \emph{the same file} changes in two places,
or changes again while an earlier version is still uploading. The vault is laid out so that never happens:

\begin{center}
\small
\begin{tabularx}{\linewidth}{@{}lX@{}}
\toprule
Path in the vault folder & Rule\\
\midrule
\file{<name>.bwvault} & The \textbf{ref file}. Identifies the vault (name, id, format version). Written once.
  Today \code{bw-hooks vault status --vault <file>} opens it; the FreeCAD panel (\eref{PDM-04}) will open it
  from File $\rightarrow$ Open or a double-click.\\
\file{objects/ab/<sha256>.FCStd} & Every checked-in file version, named by its content hash. \textbf{Written once,
  never changed}; written under a \file{.tmp} name (which iCloud does not sync) and renamed when complete; checked
  against its hash when read.\\
\file{journal/<user>@<mac>.jsonl} & Every action (create, check out, check in, submit, approve, ECR, ECO, BOM) as
  one line. \textbf{Append-only, and only that user on that Mac writes it.} The state is rebuilt by replaying the
  lines through the reference model (\eref{PDM-02}), so the lifecycle rules are identical.\\
working folder (outside the vault) & Where FreeCAD opens and saves while you work. Check-in copies the file into
  \file{objects/}; check-out copies the latest one out. Saving every few minutes never touches the vault.\\
\bottomrule
\end{tabularx}
\end{center}

If a sync service makes a copy anyway (iCloud \file{x 2.jsonl}, Dropbox \file{x (conflicted copy).jsonl},
Syncthing \file{x.sync-conflict-\ldots}), the reader reads \emph{every} journal file and keeps each event once, so a
copy neither loses nor doubles anything. A half-synced last line is skipped until the rest arrives; a journal line
whose file has not downloaded yet is reported as ``not synced yet'' rather than an error. Moving to another
provider means moving the folder.

\paragraph{Built and tested headless.} \file{addons/BruceWorx/bruceworx/pdm\_folder.py} with the command line
\code{bw-hooks vault} (actions \code{init}, \code{status}, \code{create}, \code{checkout}, \code{checkin},
\code{undo-checkout}, \code{submit}, \code{approve}, \code{reject}, \code{fetch}). The 12 lifecycle
and BOM contract tests of \eref{PDM-02} pass unchanged against a real folder, plus 17 folder tests: reopening rebuilds
identical state; a refused action writes nothing; versions are never rewritten; the released revision freezes its
file; damaged files are caught; three kinds of duplicate journal, an emptied journal, a half-written line, a missing
download, and a duplicated ref file. The OpenAPI file stays as the list of operations; the folder vault
implements them locally, and a server can implement them later.

\subsubsection*{iCloud: what we found}

\begin{description}[leftmargin=0pt]
  \item[Your Documents folder is in iCloud.] Desktop \& Documents sync is on
        (\file{\textasciitilde/Library/Mobile Documents/com\textasciitilde apple\textasciitilde CloudDocs/Documents} points
        at \file{\textasciitilde/Documents}), so this repository is synced.
  \item[Where the duplicates are.] 1\,488 numbered copies, \emph{all} in machine-written trees: build output
        (\file{build/release/share}, \file{src}) and 21 inside \file{.git} (\file{HEAD 2}, \file{index 2},
        \file{config 2}). None among the hand-edited sources. They arrived in bursts at clone and build times
        (937 in one minute on 15 Oct). The copies hold older contents: iCloud kept the version it had while the
        build was replacing thousands of files.
  \item[Single-file saves did not reproduce it.] On 19 Oct, 120 rapid rewrites of three files (save by rename,
        delete and recreate, overwrite in place) and three write passes over 3\,000 small files on this Mac produced no
        copies. The trigger is heavy, fast, machine-made churn, which is exactly what the vault never does.
  \item[Excluding a folder works on this Mac (macOS 27.2).] A folder whose name ends in \file{.nosync}, or one
        carrying the attribute \code{com.apple.fileprovider.ignore\#P}, is reported by \code{fileproviderctl} as
        \code{isExcludedFromSync = 1} and is not uploaded. Both are undocumented by Apple and could change; the
        \file{.nosync} name is the more widely used.
  \item[Storage optimisation.] With ``Optimise Mac Storage'' on, iCloud can evict files to save space; opening one
        waits for a download and fails offline. For the vault folder, turn it off or use Finder's \emph{Keep
        Downloaded} on the folder.
  \item[Sharing needs a folder in iCloud Drive itself.] Apple does not share the Desktop or Documents folders
        themselves; a shared folder has to sit in iCloud Drive. Folder sharing needs macOS 10.15.4, iOS 13.4 or
        iCloud for Windows 11.1 or later.
\end{description}

\paragraph{Recommended setup (yours to apply).}
\begin{enumerate}[nosep]
  \item Put the vault in \textbf{iCloud Drive, not Documents}, e.g.\ \file{iCloud Drive/BruceWorx Vault/}. A vault in
        Documents works for you alone, but a vault in iCloud Drive can be shared later without moving it.
  \item Keep the working folder out of sync: \file{\textasciitilde/BruceWorx Work} (outside Documents), or a
        \file{Work.nosync} folder.
  \item Turn on \emph{Keep Downloaded} for the vault folder.
  \item This repository: move it out of Documents (for example to \file{\textasciitilde/Developer}), or at least
        exclude \file{build} and \file{.pixi}. Git inside iCloud is a known way to damage a repository, and its
        \file{.git} already holds 21 duplicate files. Not changed here; it is your machine and your call.
\end{enumerate}

\paragraph{The repository, looked at again (19 Oct).} \code{git fsck} is clean, but the duplicates show the risk.
\file{.git/HEAD 2} reads \code{ref: refs/heads/.invalid} and \file{.git/config 2} has lost the branch and submodule
settings: both are the half-written state from the clone. Had iCloud settled those conflicts the other way, the
repository would have lost track of its branch and of GitHub. Excluding only \file{build} and \file{.pixi} would
leave this in place, and a branch switch or rebase rewrites thousands of source files at once, which is the same
burst that made the copies.
\begin{description}[leftmargin=0pt,nosep]
  \item[Recommended.] Repository in \file{\textasciitilde/Developer}, outside iCloud. Its off-machine copy is GitHub
        (a pushed work branch) plus Time Machine, both of which understand git. The vault, the one thing that needs
        file sync, goes in \file{iCloud Drive/BruceWorx Vault/}, which is built for it.
  \item[Not recommended.] Keeping the repository in Documents with \file{.git}, \file{build} and \file{.pixi}
        excluded. The sources would still sync, but iCloud's copy would be sources without history, and the exclusion
        relies on undocumented behaviour.
\end{description}
Moving needs a fresh \code{pixi install} and a full build, because both store absolute paths. Uncommitted work should
be pushed to a branch first, so that it exists somewhere other than this Mac.

\paragraph{Limits of this step.} One writer at a time. Journals are per user and Mac, so a second Mac or a second
person can be added later, but two people editing at the same moment are not reconciled, and the check-out lock is a
record in the journal, not a server lock. When a team needs that, the same operations move behind a server
(option A below) without changing the files.

%======================================================================
\subsection{BW-PDM-01 --- Evaluation memo (19 Oct 2026, before the decision)}\label{sec:pdm-eval}
%======================================================================

\textbf{Status.} Kept for the record. Option C was scored as a \emph{plain} shared folder with FreeCAD files
saved straight into it; the decision above is a folder vault built to avoid exactly the weaknesses scored here.

\begin{center}
\scriptsize
\setlength{\tabcolsep}{3pt}
\renewcommand{\code}[1]{\texttt{#1}}
\begin{longtable}{@{}>{\raggedright\arraybackslash}p{2.3cm}>{\raggedright\arraybackslash}p{3.4cm}>{\raggedright\arraybackslash}p{3.4cm}>{\raggedright\arraybackslash}p{3.2cm}>{\raggedright\arraybackslash}p{3.3cm}@{}}
\toprule
Evidence & A. Ops is the PLM & B. Ops + Git LFS & C. Vault on a shared folder & D. Open-source PLM\\
\midrule\endhead
\bottomrule\endfoot
Platform today & Ops = Next.js 16 on Railway (\code{bruceworks}); task API and auth are being built on branches (staging \code{/api/tasks} answers 401) & git-lfs 3.8 installed here; GitHub hosts the repos already & Folder on NAS or cloud drive & DocDokuPLM (Java), LibrePLM (Odoo; Windows-only connector), FreePDM (concept, no releases)\\
File storage cost & Railway bucket: \$0.015/GB-month, S3 operations and bucket egress free & GitHub Free/Pro: 10\,GiB storage + 10\,GiB bandwidth per month, then \$0.07/GiB storage, \$0.0875/GiB bandwidth; Team: 250\,GiB each & Disk you own & Self-hosting + upkeep\\
Locking & Server-side lease per item (contract: \code{/checkout}, \code{/checkin}) & \code{git lfs lock}: others cannot push a locked file & None real. \textbf{This repo already shows the failure:} iCloud left conflict copies (\file{BruceWorx 2.py}, \file{build 2.ninja}) instead of refusing & Built in\\
Lifecycle and ECO & Exactly the reference model (\file{bruceworx/pdm.py}) & In Ops (same model); Git only versions files & Must be built in the FreeCAD client & Configurable; fit with Y14.35 + ECR/ECO unproven\\
Multi-file assembly links & Client maps vault paths to a local workspace & Repo-relative paths work as they are & Share-relative paths & Needs a FreeCAD connector (none exists)\\
BOM and where-used & Ops (graph queries) & Ops & Client index & Built in\\
Agents / MCP & One API for FreeCAD, agents and the website & Git + Ops API & Weak & Varies\\
Main risk & Effort: we own a vault & Two systems to keep in step; CAD users must use Git (or a client that hides it) & Silent duplicates, no server truth & Connector effort, maintenance, Windows bias\\
\end{longtable}
\end{center}

\begin{center}
\small
\begin{tabular}{@{}lccccc@{}}
\toprule
Criterion (proposed weight) & A & B & C & D\\
\midrule
Data integrity and locking (25\,\%) & 5 & 4 & 2 & 4\\
Lifecycle fit, Y14.35 + ECR/ECO (20\,\%) & 5 & 3 & 3 & 3\\
Multi-file references (15\,\%) & 4 & 5 & 4 & 2\\
BOM and where-used (15\,\%) & 5 & 4 & 2 & 5\\
Build effort, higher is less work (10\,\%) & 2 & 3 & 3 & 2\\
Agent / API access (10\,\%) & 5 & 4 & 2 & 3\\
Running cost (5\,\%) & 5 & 3 & 5 & 4\\
\midrule
\textbf{Weighted score} & \textbf{4.55} & \textbf{3.80} & \textbf{2.75} & \textbf{3.35}\\
\bottomrule
\end{tabular}
\end{center}

\paragraph{Recommendation at the time.} Prototype \textbf{A} and \textbf{B} on \code{LE-02001}, then decide.
Rule out \textbf{C}: its failure mode (silent duplicates) is already visible in this repository. Keep \textbf{D}
as a fallback only, since none has a FreeCAD connector. Either A or B uses the same reference model and contract
(\eref{PDM-02}), so no work is lost whichever wins.

\paragraph{Prototype plan (not pursued after the decision).}
\begin{description}[leftmargin=0pt]
  \item[A.] A small service implementing \file{docs/pdm-openapi.yaml} over \file{bruceworx/pdm.py}, with files in
        a Railway bucket and SHA-256 content addressing. Run the contract tests against it, then check in the
        \code{LE-02001} tree and walk one ECO. Creating the bucket and service is an outward-facing change on your
        Railway account, so it waits for you.
  \item[B.] A private GitHub repository with LFS tracking for \code{*.FCStd} and locking turned on. Two clones
        play two engineers; measure lock conflicts and the clone size after ten revisions of the pedestal.
\end{description}


\subsection{Data model}

\begin{figure}[ht]
\centering
\begin{tikzpicture}[>=Stealth,
  ent/.style={draw=bwblue!70,fill=bwblue!5,rounded corners=2pt,align=left,font=\scriptsize,inner sep=4pt,text width=3.5cm},
  rel/.style={font=\tiny\color{bwgrey}}]
\node[ent] (item) at (0,0) {\textbf{Item}\\P/N (permanent)\\type: part, assembly, drawing, COTS\\description, program};
\node[ent] (rev) at (5.0,0) {\textbf{Revision}\\code: 01, 02, \ldots, A, B, \ldots\\state: WIP, Review, Released, Obsolete\\released by, date};
\node[ent] (fv) at (10.0,0) {\textbf{FileVersion}\\content SHA-256, size\\checked in by, time\\lock holder};
\node[ent] (bom) at (5.0,-2.9) {\textbf{BOM line}\\parent revision $\rightarrow$ child item\\qty, find no., configuration\\effectivity: latest released};
\node[ent] (eco) at (10.0,-2.9) {\textbf{ECO}\\affected revisions\\approvals, release};
\node[ent] (ecr) at (10.0,-5.1) {\textbf{ECR}\\problem, requester\\affected items, status};
\draw[->] (item) -- node[rel,above]{has 1..*} (rev);
\draw[->] (rev) -- node[rel,above]{has 1..*} (fv);
\draw[->] (rev) -- node[rel,right]{parent} (bom);
\draw[->] (bom.west) -| node[rel,pos=0.25,below]{child} (item.south);
\draw[->] (ecr) -- node[rel,right]{approved $\Rightarrow$ opens} (eco);
\draw[->] (eco.north west) -- node[rel,right,pos=0.35]{creates, releases} (rev.south east);
\end{tikzpicture}
\caption{PDM data model. The P/N belongs to the \emph{Item} and never changes. Revisions carry the state.
File versions are content-addressed. BOM lines point from a parent \emph{revision} to a child \emph{item},
so a released assembly picks up the child's latest released revision unless it is pinned.}
\label{fig:pdm-model}
\end{figure}

\begin{epic}{PDM-02}{Data model \& lifecycle spec}{M2}{\stProg}{(new)}
Specify Figure~\ref{fig:pdm-model} and the lifecycle in Figure~\ref{fig:lifecycle-sm} as an API contract
(OpenAPI) that any storage option implements. A drawing is its own \emph{Item}, numbered like its part, whose
revision moves \emph{with} the part's: one ECO releases both.
\textbf{Drafted 19 Oct 2026:} reference model \file{addons/BruceWorx/bruceworx/pdm.py} (revision sequences,
lifecycle, ECR/ECO, drawing coupling, BOM explosion, where-used, cycle rejection) with 18 contract tests in
\file{tests/test\_pdm.py}, and the API contract \file{addons/BruceWorx/docs/pdm-openapi.yaml} (17 endpoints,
including check-in/out and locks).
\emph{Acceptance:} spec approved; every transition in Figure~\ref{fig:lifecycle-sm} has an endpoint and a test.
\end{epic}

\subsection{Lifecycle and revisions}\label{sec:lifecycle}

\begin{figure}[ht]
\centering
\begin{tikzpicture}[>=Stealth,x=3.6cm,y=-2.2cm,
  s/.style={draw,rounded corners=3pt,minimum width=2.3cm,minimum height=0.85cm,align=center,font=\small},
  wip/.style={s,draw=bworange,fill=bworange!8},rv/.style={s,draw=bwblue,fill=bwblue!8},
  rl/.style={s,draw=bwgreen,fill=bwgreen!10},ob/.style={s,draw=black!40,fill=black!5},
  ch/.style={s,draw=bwred,fill=bwred!6},
  t/.style={font=\scriptsize,align=center}]
\node[wip] (wip) at (0,0) {WIP\\\scriptsize rev 01, 02, \ldots};
\node[rv]  (rv)  at (1,0) {In Review};
\node[rl]  (rl)  at (2,0) {Released\\\scriptsize rev A};
\node[ob]  (ob)  at (3,0) {Obsolete};
\node[ch]  (ecr) at (2.9,1) {ECR};
\node[ch]  (eco) at (1.9,1) {ECO\\\scriptsize work rev B};
\node[rv]  (rv2) at (0.9,1) {In Review};
\draw[->] (wip) -- node[t,below]{submit} (rv);
\draw[->] (rv.north) to[bend right=35] node[t,above]{reject: next numeric rev} (wip.north);
\draw[->] (rv) -- node[t,above]{approve} (rl);
\draw[->] (rl) -- node[t,above]{retire} (ob);
\draw[->] ([xshift=0.6cm]rl.south) -- ++(0,0.22) -| node[t,pos=0.25,above]{problem reported} (ecr.north);
\draw[->] (ecr) -- node[t,above]{approved} (eco);
\draw[->] (eco) -- node[t,above]{submit} (rv2);
\draw[->] (rv2.west) -- ++(-0.1,0) |- node[t,pos=0.72,below]{approve: rev B released, rev A superseded} ($([xshift=-0.5cm]rl.south)+(0,-0.45cm)$) -- ([xshift=-0.5cm]rl.south);
\end{tikzpicture}
\caption{Lifecycle. Pre-release revisions are numeric (ASME Y14.35 practice) and every review rejection
increments them. The first approval releases rev A. After that, letters change only through an approved ECR
$\rightarrow$ ECO, and released files are read-only. ``Superseded'' is a flag on the old letter revision, not a state.}
\label{fig:lifecycle-sm}
\end{figure}

\paragraph{Theory: revision sequences.} Released revisions use letters with I, O, Q, S, X and Z skipped,
as ASME Y14.35 recommends, because they read as digits. Writing the 20 usable letters as an alphabet
$\Sigma$, the sequence is base-20 with no zero digit: A, B, \ldots, Y, AA, AB, \ldots. The successor
function increments the last letter and carries. Pre-release revisions are zero-padded integers. The drawing
prints \code{P/N-REV}, for example \code{LE-01-01001-001-02} or \code{LE-01-01001-001-A}, while the P/N
stored on the item never includes the revision (decided).

\paragraph{Theory: why pessimistic locking.} \code{.FCStd} files are zip archives of XML and B-rep data. Two
concurrent edits cannot be merged, so optimistic concurrency, where both editors save and a merge resolves
it, would always end with one person's work discarded. With 2--5 editors, \emph{pessimistic} locking (check-out
grants an exclusive lock; check-in releases it) costs little waiting and avoids lost work. Locks carry a lease,
and an administrator can break a stale one. A file version is identified by its SHA-256, so an unchanged
check-in is a no-op and duplicate content is stored once.

\begin{epic}{PDM-03}{Numbering service}{M2}{\stProg}{FCB-P2-03/04/05}
The local-JSON numbering source and macro are merged (\#5). The format is
\code{\{PROG\}-\{DT\}-\{PPPPP\}-\{SSS\}-\{REV\}}. Under the decision, \code{\{REV\}} moves out of the P/N into
revision metadata, and label and file names drop it. Then the Ops API reader (stubbed as
\code{OpsApiNumberSource}) allocates numbers centrally, with a bearer token from \code{BW\_OPS\_TOKEN}.
\emph{Acceptance:} P/Ns come from Ops; existing LE numbers migrate with only the revision suffix removed.
\end{epic}

\begin{epic}{PDM-04}{FreeCAD vault client}{M2}{\stTodo}{(new)}
A Vault panel and tree overlays: \emph{checked in}, \emph{checked out by me}, \emph{locked by someone
else}, \emph{newer version available}. Commands: check out, check in (with comment), undo check-out, get
latest, and get a specific revision. Opening a locked file opens it read-only. Saving a file you have not
checked out offers to check it out first. Assembly links resolve through the vault, so \code{LE-02001}
opens identically on every machine.
Works over the folder vault (\S\ref{sec:pdm-decision}); the command-line \code{vault} commands already
drive it.
\emph{Acceptance:} a check-out/check-in round trip of \code{LE-02001} from an iCloud Drive vault inside FreeCAD;
the status icons are always correct. Locking between several users is a later step.
\end{epic}

\begin{epic}{PDM-05}{Lifecycle \& change control}{M2}{\stTodo}{(new)}
Implements Figure~\ref{fig:lifecycle-sm}, with approvals in Ops. Release freezes files read-only, publishes
drawings (\eref{DRW-08}), and builds the release package (\eref{MFG-04}). The ECR captures the problem;
the ECO groups affected items, opens their next letter revision in WIP, and releases them together.
\emph{Acceptance:} rev A $\rightarrow$ ECR $\rightarrow$ ECO $\rightarrow$ rev B round trip on \code{LE-01001}
and its drawing.
\end{epic}

\begin{epic}{PDM-06}{BOM sync \& where-used}{M2}{\stTodo}{(new)}
On check-in of an assembly, push its BOM (\eref{ASM-04}) to Ops as BOM lines. Where-used is the reverse
traversal from \S\ref{sec:asm}. Released assemblies resolve children to the latest \emph{released}
revision. A WIP assembly may also use WIP children, and is flagged when it does.
\emph{Acceptance:} where-used of a bolt lists every assembly revision that uses it.
\end{epic}

\begin{epic}{PDM-07}{Pilot: LE-02001 \& CAD register}{M2}{\stTodo}{FCB-P2-15/16/17/18, P2-21}
The first real release: build the pedestal tree (\code{LE-01001} $\rightarrow$ \code{LE-01010}
$\rightarrow$ \code{LE-01009} $\rightarrow$ \code{LE-01011} $\rightarrow$ \code{LE-01002}), link the CAD
register in the BruceWorx app, run the design-freeze checklist and validation macro, and release rev A with
drawings and package.
\emph{Acceptance:} \code{LE-02001} rev A released entirely from FreeCAD.
\end{epic}

\begin{epic}{PDM-08}{Migration from Fusion}{M6}{\stTodo}{FC-T01, FCB-P2-22}
Bulk-import Fusion exports (STEP, plus drawings as PDF) into the vault with P/N and revision mapping,
re-establish assembly structure, and record provenance. The FC-T01 switch criteria become the checklist for
M6.
\emph{Acceptance:} every LE part imported, numbered and released; FC-T01 criteria met.
\end{epic}
